\section{Motivation: Every Token Has a Cost}
\label{sec:motivation}

Large language models pay for every token they read. That cost shows up
in the energy and water to read the prompt and generate
the reply (the prefill and decoding passes); the latency the user waits
through; and the cost the user pays, whether per million tokens
through an API or amortised over on-premises hardware. All three scale
with the number of tokens, and recent but consistent peer-reviewed work
bears this out. Luccioni \emph{et~al.}'s 2024 \emph{Power Hungry
Processing} study reports per-inference energy across a range of open
models, about $0.1$~Watt-hours for a 6--7B-class
model~\cite{luccioni2024powerhungry}; Samsi
\emph{et~al.}~\cite{samsi2023words} show that inference energy scales
with the number of tokens processed; and Pope
\emph{et~al.}~\cite{pope2023efficiently} quantify how serving cost grows
with sequence length once the prefill dominates. These inference-time
studies extend the training-time carbon-and-cost accounting that
Patterson \emph{et~al.}~\cite{patterson2021carbon} established. And when
the workload is a multi-agent system that issues thousands of
internal inference calls per user request, the token costs
compound.

Yet a growing fraction of the tokens a deployed LLM reads do not change
the answer. They are conversational filler (``Hello'', ``Thanks''),
repeated system prompts that scaffold every call, restated history that
re-anchors a long conversation, boilerplate disclaimers, and framing that could be skipped. Because cost is tied to the tokens
processed~\cite{luccioni2024powerhungry, samsi2023words,
pope2023efficiently}, each such token still turns into energy, latency,
and operating expense whether or not it has an impact on the outcome.

As a consistent industry anecdote, OpenAI's compute bill for processing
user politeness alone was informally estimated at ``tens of millions of
dollars'' a year in 2025~\cite{altman2025pleasethankyou}.

Smart context compression is the natural solution. Per-token compute and cost
scale roughly linearly in sequence length once the prefill is past
the GPU memory-bound regime, so halving the input context roughly
halves the cost. The single-agent prompt-compression literature has
matured fast over the past three
years~\cite{jiang2023llmlingua, jiang2024longllmlingua, pan2024llmlingua2,
mu2023gist, chevalier2023autocompressor, ge2024icae}. This thesis
measures what that literature has left unanswered: whether compression
that preserves single-fragment quality (the single-agent
question-answering metric it optimises) also preserves information spread across several fragments into one correct answer.

\section{The Multi-Fragment Workflow}
\label{sec:multi_fragment}

This thesis studies the \emph{multi-fragment LLM workflow}: a task
whose answer cannot be read out of any single fragment but instead
requires combining information across two or more fragments, each of
which may be compressed on its own before the planner sees it. Three
examples recur: cross-document fact aggregation (sum eight numbers
reported by eight institutional systems), constraint-satisfaction
planning (assign $N$ sub-tasks to $K$ capacity-bounded workers from
specifications spread across fragments), and multi-step retrieval
(chain references through a sequence of documents to a leaf answer);
the C1 benchmark of \charef{implementation} builds $50$ instances of
each.  The FCG group at the
University of Oulu's flagship use case,
\emph{Vignette~3.7 -- Period-end project reporting}~\cite{fcgusecase2026},
gathers reporting data across eight university systems via roughly
eight specialised agents; the internal FCG analysis estimates an
$80\%$ cloud-payload reduction at about EUR~$2.15$ per report.
TalentAdore, the industry partner of this thesis, runs production
workloads whose operating expense is dominated by token cost.
Anthropic's report on their multi-agent research system finds token
usage explains about $80\%$ of the performance variance they
observe~\cite{anthropic2025multiagent}, and the companion
context-management report describes an $84\%$ token reduction on a
100-turn web-search evaluation~\cite{anthropic2025contextediting}.
These are industry observations on industry workflows; we reshaped the problem into a controlled cross-compressor,
cross-task, cross-architecture measurement that turns that urgency
into a technical question.

The technical question deliberately excludes the multi-round agent
negotiation that the ``multi-agent'' framing of the FCG and Anthropic
systems implies. The empirical evaluation uses one of two planners: a
deterministic regex information-extraction solver, or a single LLM
call with all (compressed) fragments visible. An AutoGen v0.4
multi-round backend ships in the source tree
~\cite{wu2023autogen} but is not used
in any reported result. The reason is attribution: round-to-round LLM
variance dominates the compression signal at the C1 ratios swept, so
cleanly attributing a coordination drop to compression rather than to
agent variance requires isolating the compressor on the critical
path, which a single-LLM-call or deterministic planner does. The
memory bus is \emph{designed for} multi-agent integration; the
experimental scope is multi-fragment task solvability, restated and
calibrated in \charef{implementation}~Section~\ref{sec:scope}.

\paragraph{Operational definition of ``coordination success.''}
Because the title and the coinage \emph{coordination cliff} invite a
multi-agent reading that the experiments do not test, the term is fixed
here. \textbf{Coordination success is the binary indicator that the
planner recovers the correct answer to a multi-fragment task whose
solution cannot be read out of any single fragment alone}: a
measurement of \emph{multi-fragment solvability under compression}, not
of multi-round agent negotiation, message-passing efficiency, or role
allocation under disagreement. Every later use of ``coordination''
refers to this quantity.

\section{Problem Statement}
\label{sec:problem}

\emph{How does context compression affect multi-fragment
coordination quality across compressors, task structures, and
planner scales, and when and how can a memory-bus service
exploit that understanding to compress safely, predictably, and
privately?}

The question subdivides into four sub-questions, one per
contribution (C1--C4), each answered end-to-end by an artefact
and operationalised by one or more of the falsifiable hypotheses
(H1--H6) that structure \charef{experiments}.

\subsection{Research questions and their mapping to hypotheses}
\label{sec:rq_mapping}

The four research questions and their hypothesis mappings are:

\begin{description}\setlength\itemsep{2pt}
\item[RQ1 (metric).] Does a coordination-quality metric on a
multi-fragment workflow behave differently from single-agent
question-answering accuracy? (operationalised by \textbf{H1}).
\item[RQ2 (characterisation).] Does compression induce a
coordination cliff, and what determines its position?
(operationalised by \textbf{H2} (existence and sharpness),
\textbf{H5}/\textbf{Corollary~1} (invariance to planner scale),
and \textbf{H6}/\textbf{Corollary~2} (scaling with task
information density)).
\item[RQ3 (deployment).] How should compression be combined with
retrieval, and what is the cost trade-off across pipeline
architectures? (operationalised by \textbf{H3}).
\item[RQ4 (governance).] Can compression be measured as a privacy
lever and exposed as a deployment-time enforcement mechanism by a
memory-bus service? (operationalised by \textbf{H4} and the
memory-bus integration of \charef{experiments}).
\end{description}

\section{Contributions}
\label{sec:contributions}

\paragraph{Headline.}
The thesis's single most-citable result is the H1 finding (the
methodological backbone of C2): on the multi-fragment benchmark
shipped as C1, the change in single-fragment source-retention F1
under compression does not positively predict the change in
multi-fragment coordination success for \emph{any} of the four
compressors tested. The single-agent benchmarks that dominate the
prompt-compression literature (LLMLingua-2, LongLLMLingua, LongBench,
RULER) therefore systematically mis-rank compressors for
multi-fragment deployment. The thesis turns that finding into its
principal engineering deliverable: a \emph{compression-aware memory bus} (C4)
whose compression layer is tuned by the cliff measurements rather
than by guesswork, built to cut inference cost in production.

\paragraph{C1: A reproducible multi-fragment coordination
benchmark.}
\emph{The C1 benchmark} ships $150$ instances across three
workload families (cross-document fact aggregation,
constraint-satisfaction planning, multi-step retrieval), with
synthetic provenance and access-control tag distributions, four
coordination-quality metrics, and single-command regeneration from a
fixed seed. It is the substrate on which the cliff characterisation,
the model-scaling result, and the disclosure measurement all run;
releasing it lets other groups extend the characterisation to new
compressors and planner regimes.

\paragraph{C2: Empirical characterisation of the coordination
cliff, a compounding-error model that explains it, and two
structural corollaries.}
\emph{The cliff exists and is sharp; its position shows no
detectable shift across the heterogeneous planners we could test
within a defined calibrated regime; and its shape varies
systematically with task information density.} The empirical work
covers four training-free compressors (LLMLingua-2~\cite{pan2024llmlingua2},
a Phi-3-Mini extractive span selector, an instruction-aware filter,
and a truncation baseline) at ten ratios across three families and
three local architecture families (Qwen-2.5-1.5B, Phi-3-Mini-3.8B,
Llama-3.1-8B), with a frontier consistency arm (Qwen-2.5-72B, DeepSeek~V4~Pro). The
\emph{compounding-error model} derives the cliff's threshold
structure from per-round \textbf{critical-token-recall} against a
per-family threshold $\theta_q$, recovering its position only to
first order: a mechanism, not a tight predictor.
\textbf{Corollary~1 (Ceiling--Cliff Separation)} and
\textbf{Corollary~2 (Information-Density Scaling)}, sharpened from
H5 and H6, capture the no-shift and shape-scaling results, the
latter via an AUC-based per-task $\theta_\text{info}$ (distinct from
$\theta_q$) validated on HotpotQA and MultiHop-RAG.

\paragraph{C3 (secondary): a catalogue of three
retrieval-augmented-generation pipeline placements.}
Three placements (compress-first, retrieve-first, joint
relevance-conditional routing) are compared under matched storage
and accuracy-bounded regimes. The pre-specified sign-flip is
\emph{not} observed; the narrowed finding is that compress-first
beats retrieve-first on the single stack tested, not reproducing the
LongLLMLingua~\cite{jiang2024longllmlingua} post-retrieval
preference. As a single-stack non-reproduction orthogonal to the
cliff and the bus, its full treatment is deferred to
Appendix~\ref{appendix:rag}.

\paragraph{C4: A compression-aware memory bus, tuned by the cliff
measurements.}
The principal engineering contribution is a memory bus that sits
between the data sources and the planner, compresses every fragment
in transit, and turns the token savings into lower inference cost: a
FastAPI service with a tamper-evident SQLite audit log, a policy
middleware enforcing a five-tier classification lattice over a
$64$-bit access-control mask, and a FAISS-CPU vector store, behind a
single \texttt{Compressor} protocol. What makes it more than
plumbing is that its compression layer is \emph{driven by the
findings}: a cliff-aware wrapper (\textbf{CAAC}) takes the
per-family threshold $\theta_q$ and backs off, per fragment, to the
most aggressive ratio that keeps critical-token-recall on the safe
side of the cliff; and a summary-level inference-disclosure metric
(H4, Section~\ref{sec:h4_disclosure}) caps how much a downstream
reader can recover a protected fact from a compressed
\textsc{confidential} fragment, enforced at write time. The research
content is the disclosure metric and the CAAC construction; the
service operationalises them. Single-host throughput is benchmarked
(Section~\ref{sec:bus_bench}); distributed scaling is follow-on work.
The implementation is scheduled for deployment in the production
workflows of the industry partner, TalentAdore, whose operating
expense is dominated by token cost AI operations.

\subsection{Reading guide and contribution hierarchy}
\label{sec:reading_guide}

Table~\ref{tab:contribution_map} is the map of the whole
thesis: it ties each research question to the artefact that answers
it.

\begin{table}[htbp]
\centering
\caption{Contribution map.}
\label{tab:contribution_map}
\small
\begin{tabularx}{\linewidth}{>{\raggedright\arraybackslash}p{1.2cm}
                              >{\raggedright\arraybackslash}p{3.2cm}
                              >{\raggedright\arraybackslash}p{1.6cm}
                              >{\raggedright\arraybackslash}X}
\hline
RQ / C & Plain-English question & Hyp. & Verdict (one line) \\
\hline
RQ1\newline C1--C2 & Does keeping a single fragment answerable also keep several fragments \emph{combinable}? & H1 & \textbf{No}: single-fragment F1 does not predict coordination success (\textsc{supported}). \\
RQ2\newline C2 & Is there a sharp failure point as compression rises, and what fixes its position? & H2; Cor.~1; Cor.~2 & A sharp cliff exists (\textsc{supported}, $11/12$ cells); its position is set by compressor$+$task, not planner scale (no shift detected); its shape scales with information density (\textsc{supported}, $n{=}3$). \\
RQ3\newline C3 & Where should compression sit relative to retrieval? & H3 & Compress-first dominates on the tested stack; the sign-flip is \emph{not} observed (\emph{secondary}; Appendix~\ref{appendix:rag}). \\
RQ4\newline C4 & Can the cliff findings drive a deployable memory bus that compresses safely and cheaply? & H4, CAAC & The bus compresses to the cliff-safe ratio (CAAC) and caps protected-fact disclosure (\textsc{supported}, $3/4$ compressors; destruction-driven). \\
\hline
\end{tabularx}
\end{table}

The four contributions are not equally weighted: \textbf{H1} is the
central, load-bearing result and the coordination cliff with its
compounding-error model (C2) is the mechanism behind it; C1 is the
enabling instrument; and the RAG catalogue (C3), the
inference-disclosure metric and memory bus (C4), and CAAC are
deliberately secondary, each operationalising the central finding in
one deployment setting.

\section{Thesis Outline}
\label{sec:outline}

\charef{relatedwork} surveys the prior work the four contributions
build on: the transformer cost of context, the scaling era, the
economics of inference, the training-free compression literature,
multi-agent systems and agentic memory, retrieval-augmented
generation and long-context benchmarks, and privacy in compressed
retrieval, closing with the gap this thesis fills.

\charef{implementation} presents the artefacts the evaluation runs
on: the three-layer memory-bus architecture, the compressor
framework, the C1 benchmark and its scoring pipeline, the
critical-token-recall metric, and the compression cache.

\charef{experiments} presents the empirical work: H1 (information
preservation does not transferably predict coordination), H2 (a
sharp cliff across $12$ compressor-family cells), the
compounding-error model that explains it, Corollary~1 (no cliff-%
position shift across planner scale and architecture within a
calibrated regime), H3 (RAG placement), H4 (inference disclosure),
and Corollary~2 (the cliff shape scales with information density).

\charef{summary} synthesises the findings, presents CAAC as the
constructive realisation of the model, and gives the limitations,
significance, and future work.


All artefacts (manuscript source, code, configuration files,
results CSVs, and the \texttt{docker-compose.yml} that brings the
memory-bus service up locally) are in the reference repository at
\url{https://github.com/Abdullah12has/Distributed-memory-bus-for-multi-agent-coordination}.
Appendix~\ref{appendix} describes the reproducibility package and the reproduction.
